\section{Bedarfsdeckung}

\subsection{Deckung des Energiebedarfs}
\subsubsection{Begriffe}
    {\small 
        \begin{itemize}
            \item \textbf{Bruttostromverbrauch: } Gesamtstromverbrauch + Eigenbedarf Kraftwerke
            \item \textbf{Gesamtstromverbrauch: } Nettostromverbrauch + Netzverluste + Pumpspeicherverbrauch
            \item \textbf{Nettostromverbrauch: } Stromverbrauch der Endkunden
        \end{itemize}
    }

\subsubsection{Zusammensetzung Bedarf}
    {\small Zusammensetzung des Strombedarfs in Deutschland (2025) in TWh: \\ 
    }
    \begin{table}[H]
        \begin{tabular}{l r}
            \textbf{Verbraucher} & \textbf{Bedarf in TWh} \\
            \toprule
            Industrie & 190 \\
            Haushalte & 130 \\
            Gewerbe, Handel, Dienstleistungen & 130 \\
            Verkehr & 20 \\
            \midrule
            \textbf{Gesamtverbrauch} & 470 \\
            \bottomrule
            Netzverluste & 30 \\
            Nettostromerzeugung & 480 \\
        \end{tabular}
    \end{table}

\subsubsection{Zusammensetzung Erzeugung}
    Gesamterzeugung: 440 TWh (2025) 
    \begin{table}[H]
        \begin{tabular}{l r}
            \toprule
            \textbf{Erzeuger} & \textbf{Erzeugungsenergie in TWh} \\
            \midrule
            Erdgas & 13,5\% \\
            Steinkohle & 6,4\% \\
            Braunkohle & 16\% \\
            Windenergie & 31,5\% \\
            Photovoltaik & 16,9\% \\
            Wasserkraft & 4,2\% \\
            Biomasse & 8,6\% \\
            sonstige & 2,9\% \\
            \bottomrule
        \end{tabular}
    \end{table}

\subsection{Deckung des Leistungsbedarfs}
\subsubsection{Leistungsgleichgewicht}
    {\small
        \underline{Stationärer Betrieb: }
        \begin{align*}
            W_{rot} = \frac{1}{2} \cdot J \cdot \omega_{mech}^2 \quad ; \quad
            W_{m_{zu}} = W_{el_{ab}} \quad ; \quad 
            P_{m_{zu}} = P_{el_{ab}}
        \end{align*}
        \underline{Kopplung zwischen Leistung und Drehzahl}
        \begin{gather*}
            \omega_{el} = p \cdot \omega_{m} \quad ; \quad
            2 \pi f = p \cdot 2 \pi \cdot n \rightarrow f = p \cdot n 
        \end{gather*}

        n: Drehzahl \quad p: Polpaarzahl \quad f: Frequenz 

        Leistungsbedarf steigt \textrightarrow Drehzahl sinkt \textrightarrow Frequenz sinkt \textrightarrow Erzeugung muss erhöht werden

        Steigt die Frequenz zu weit, werden Erzeuger abgeschaltet. \\
        Sinkt die Frequenz zu weit, werden Verbraucher abgeworfen.
    }

\subsubsection{Belastungsdiagramm}
    {\small
        \begin{itemize}
        \item Die Nennbetriebsdauer \textbf{\(T_n\)} in h (=Beobachtungsdauer) gibt Art des Belastungsdiagramms an. 

        \item \textbf{Kenngrößen: }
            \begin{itemize}
                \item Nennleistung \(P_n\) und Ausnutzungsdauer \(T_a\) (maximal mögliche Leistung)
                \item Höchstlast \(P_{max}\) und Benutzungsdauer \(T_m\) (maximal genutzte Leistung)
                \item mittlere Leistung \(P_{mittel}\) und Nennbetriebsdauer \(T_n\) (Gesamtzeit)
            \end{itemize}

        \item \textbf{Elektrische Energie: }
            \begin{align*}
                W = \int_{0}^{T_n} P(t) \, dt = P_n \cdot T_a = P_{max} \cdot T_m = P_{mittel} \cdot T_n
            \end{align*}
        
        \item \textbf{Typische Jahresnutzungsdauer \(T_m\): }
            \begin{table}[H]
                \centering
                \small
                \begin{tabular}{l l}
                    \toprule
                    \textbf{Verbraucher} & \\
                    \midrule
                    HöS-Netz: Starklast & 7660 h (365 x 21 h) \\
                    HöS-Netz: Schwachlast & 7120 h (365 x 19,5 h) \\
                    öffentliche EV & 6200 h \\
                    Landwirtschaft & 150 bis 250 h \\
                    Haushalt & 500 bis 1300 h \\
                    chemische Industrie & 6000 bis 8000 h \\
                    \toprule
                    \textbf{Erzeuger} & \\
                    \midrule
                    Kernkraftwerke & 8500 h \\
                    Laufwasserkraftwerke & 6500 bis 7500 h \\
                    Offshore-Windparks & 3000 h \\
                    Onshore-Windparks & 2000 h \\
                    Photovoltaik & 800 bis 1000 h \\
                    \bottomrule
                \end{tabular}
            \end{table}

        \item \textbf{Kategorisierung:} 
            \begin{table}[H]
                \centering 
                \small
                \begin{tabular}{l r r}
                    \toprule
                    Kategorie & Jahresdauer & Tagesdauer \\
                    \midrule
                    Spitzenlast & < 2000 h & < 5 h\\
                    Mittellast & 2000 h \textasciitilde \ 5000 h & 5 h \textasciitilde \ 14 h\\
                    Grundlast & > 5000 h & > 14 h\\
                    \bottomrule
                \end{tabular}
            \end{table}

            \begin{table}[H]
                \centering
                \small
                \begin{tabularx}{\columnwidth}{l X}
                    \toprule
                    Kategorie & Beschreibung \\
                    \midrule
                    Grundlast & dauerhaft vorhandene Residuallast \\
                    Regellast & zu regelnde (Residual-) Last über der Grundlast\\
                    \bottomrule
                \end{tabularx}
            \end{table}
            Residuallast = Gesamtlast - Erneuerbare Erzeugung

        \end{itemize}
    }